% Part: model-theory
% Chapter: basics
% Section: theory-of-m

\documentclass[../../../include/open-logic-section]{subfiles}

\begin{document}

\olfileid{mod}{bas}{thm}

\section{نظرية \printtoken{S}{structure}}

لكل !!{structure}~$\Struct{M}$ !!{sentence}s تصدق فيها
وأخرى تكذب. فنجمع جميع الـ!!{sentence}s الصادقة ونسمي مجموعها
\emph{نظريتها}. ويصدق عليها اسم النظرية على الحقيقة، لأن كل
ما يلزم عنها يصدق في جميع نماذجها، ولا سيما~$\Struct{M}$.

\begin{defn}
  لتكن !!a{structure}~$\Struct M$. \emph{نظرية} $\Struct{M}$ هي مجموعة
  الـ!!{sentence}s~$\Theory{M}$ الصادقة في $\Struct{M}$، أي
  $\Theory{M} = \Setabs{!A}{\Sat{M}{!A}}$.
\end{defn}

وقد نتوسع في لفظ «نظرية» فنطلقه، على غير الاصطلاح الدقيق، على مجموعات
من الـ!!{sentence}s ذات تأويل مقصود، من غير اشتراط انغلاقها استنباطيًا.

\begin{prop}
لكل $\Struct{M}$، تكون $\Theory{M}$ تامة.
\end{prop}

\begin{proof}
لنأخذ أي !!{sentence}~$!A$. فلا بد من $\Sat{M}{!A}$ أو
$\Sat{M}{\lnot !A}$؛ فتكون $!A \in \Theory{M}$ أو
$\lnot !A \in \Theory{M}$، وهو المطلوب.
\end{proof}

\begin{prop}\ollabel{prop:equiv}
  إذا كان $\Struct{N} \models !A$ لكل $!A \in \Theory{M}$، فإن
  $\Struct{M} \elemequiv \Struct{N}$.
\end{prop}

\begin{proof}
من الفرض $\Sat{N}{!A}$ لكل $!A \in \Theory{M}$، ولذلك
$\Theory{M} \subseteq \Theory{N}$. ولإثبات الاحتواء الآخر،
إذا كان $\Sat{N}{!A}$، امتنع $\Sat{N}{\lnot !A}$، فلزم
$\lnot !A \notin \Theory{M}$. وبما أن $\Theory{M}$ تامة،
وجب $!A \in \Theory{M}$. فثبت
$\Theory{N} \subseteq \Theory{M}$، ومن ثم
$\Struct{M} \elemequiv \Struct{N}$.
\end{proof}

\begin{rem}\ollabel{remark:R}
  لنأخذ $\Struct{R} = \langle\Real, <\rangle$، وهي الـ!!{structure}
  التي مجالها الأعداد الحقيقية~$\Real$، ولغتها الـ!!{language}
  ذات !!{predicate} واحد ثنائي يؤول بالعلاقة~$<$ على الحقيقية.
  لا ريب أن $\Struct{R}$~!!{nonenumerable}. لكن $\Theory{R}$
  متسقة، فلها بمبرهنة لوفنهايم--سكولم نموذج !!{enumerable}،
  نسمّيه $\Struct{S}$. ويلزم من \olref{prop:equiv} أن
  $\Struct{R} \equiv \Struct{S}$. ولما كانت $\Struct{R}$
  و$\Struct{S}$ غير متشاكلتين، بطل عكس \olref[iso]{thm:isom}
  على العموم.
\end{rem}

\end{document}
